\chapter{Linking and Loading}\label{ch:linking-loading}

Compilation ends with object files, but programs run as processes.
Linking and loading close that gap by combining objects, resolving names, assigning addresses, and preparing a runnable image.
The linker merges sections, patches relocations, and exports a symbol table.
The loader maps segments into memory, applies final fixups, and transfers control to the entry point.
Both must be exact: one wrong address computation crashes the program before it starts.
Chapters~\ref{ch:assembly-language}--\ref{ch:lua-assembler} develop the assembler side of this boundary.
Chapters~\ref{ch:link-symbols}--\ref{ch:lua-linker} expand the linking stages below into detailed algorithms and a tested Lua-mapped ELF64 linker.

\section{What the Linker Combines}\label{sec:ch29-linker-job}

A program typically spans many translation units plus libraries.
Each object file contributes named sections such as text, read-only data, writable data, uninitialized storage, debug information, and unwind tables.
\term{Static linking} merges those contributions into one executable or shared object with all internal addresses resolved.
\term{Symbol resolution} matches each undefined reference to exactly one definition under the scoping rules.
\term{Relocation} patches code and data bytes whose final addresses were unknown at compile time.
\term{Section merging} concatenates same-named input sections, discards dead ones when asked, and assigns the result to output segments with permissions.
A \term{segment} groups sections with the same memory permissions for the loader, while a section is the finer linker-level unit.

Consider a call from \code{main.o} to \code{helper} defined in \code{util.o}, plus a global counter incremented by both.
At compile time the call target is a placeholder and the counter address is a placeholder.
The symbol tables record an undefined \code{helper} in the first object, a defined \code{helper} in the second, and a common or defined counter in one of them.
Relocation records name each placeholder with its symbol, offset, and computation.
The linker picks the layout, writes the final addresses into the placeholders, and reports duplicates or missing definitions.
The loader later maps the segments and, for position-independent or dynamic builds, applies the remaining fixups.
\begin{irlisting}
; relocations in main.o (schematic)
text+0x10: call helper + addend -4  ; direct call to fill
text+0x20: load counter + GOT[tls?] ; address to fill
symtab: UND helper, UND counter
; definitions in util.o (schematic)
symtab: DEF helper at text+0x0, DEF counter at data+0x0
\end{irlisting}
Nothing here is executable until the placeholders are patched.
That patching is the linker entire visible effect on correct builds, plus layout, ordering, and dead-stripping choices that affect size and locality.
Figure~\ref{fig:ch29-merge} shows section merging from two objects into one image.
Table~\ref{tab:ch29-reloc-kinds} lists the relocation computations a linker applies most often.

\begin{figure}[htbp]
\centering
\begin{tikzpicture}[node distance=0.7cm and 1.4cm]
\node[basic block] (t1) {main.o .text};
\node[basic block, below=of t1] (d1) {main.o .data};
\node[basic block, right=of t1] (t2) {util.o .text};
\node[basic block, right=of d1] (d2) {util.o .data};
\node[basic block, right=of t2] (out) {image: text + data};
\draw[cfg edge] (t1) -- (out);
\draw[cfg edge] (t2) -- (out);
\draw[cfg edge] (d1) -- (out);
\draw[cfg edge] (d2) -- (out);
\end{tikzpicture}
\caption{The linker concatenates matching input sections into output segments, resolves symbols across objects, and patches each relocation against the final layout.}
\label{fig:ch29-merge}
\end{figure}

\begin{table}[htbp]
\caption{Common relocation computations and when they apply.}
\label{tab:ch29-reloc-kinds}
\begin{tabular}{lll}
\toprule
Computation & Meaning & Typical use \\
\midrule
Direct 32 or 64 & Symbol address plus addend & Absolute pointers \\
PC-relative 32 & Symbol minus place plus addend & Calls and jumps \\
GOT-relative & Table offset of symbol & Shared-library data \\
PLT call & Stub address for function & Lazy binding \\
TLS variant & Thread-block offset & Thread-local data \\
\bottomrule
\end{tabular}
\end{table}

Duplicate and missing symbols are the two resolution outcomes every programmer meets.
Strong definitions collide: two objects defining the same global name is an error in most models, with special rules for weak symbols, linkonce sections, and common-block merging inherited from older toolchains.
Missing definitions fail the link unless the symbol is explicitly allowed to stay undefined for runtime resolution.
Archives (static libraries) refine resolution by contributing only the members that satisfy some outstanding undefined name, which makes link order significant.
Shared objects refine it further by deferring binding to load time through tables rather than patching every use site directly.
\begin{ccnote}[Link order is semantics]
With archives, the linker scans left to right and extracts only what resolves a known undefined name.
Reordering libraries can flip a build from success to failure without changing any source.
Prefer explicit dependency lists over hopeful ordering.
\end{ccnote}

\section{Relocation and Address Assignment}\label{sec:ch29-algorithm}

Algorithm~\ref{alg:ch29-link} states the static-link skeleton: collect symbols, lay out sections, then patch relocations.
It assumes well-formed objects with consistent section flags, a link script or default layout, and a symbol table where each global has compatible type and size.
It produces an executable or relocatable image plus diagnostics for duplicates and undefined names.
Real linkers add garbage collection of unused sections, identical-code folding, ordering for locality, and relaxation that rewrites instruction sequences once distances are known.

\begin{algorithm}[htbp]
\caption{Static link: resolve, lay out, and relocate}\label{alg:ch29-link}
\begin{pseudocode}
\Require object set $O$, library set $L$, layout rules $D$
\Ensure linked image or diagnostic
\State build the symbol table: resolve definitions across $O$
\State extract needed archive members under the linker's search-order rules; update definitions and unresolved names until that search reaches closure
\State merge input sections into output sections per $D$
\State assign addresses honoring alignment and permissions
\ForAll{relocation $r$ at place $p$ naming symbol $s$}
\State compute $value(s, p)$ per the relocation kind
\State check overflow and permission for $r$
\State write the patched bytes at $p$
\EndFor
\State emit segments, symbol tables, and debug mappings
\end{pseudocode}
\end{algorithm}

Each relocation computation is a small integer formula over the symbol address $S$, the place address $P$, the addend $A$, and sometimes a table base $G$.
An extracted archive member can introduce new undefined symbols, so extraction is not necessarily a single pass over the original unresolved set.
Library ordering and explicitly grouped rescans are part of the linker's contract; a mathematical set of libraries alone does not encode those rules.
Direct absolute relocation writes $S + A$ truncated to the field width with an overflow check.
PC-relative relocation writes $S + A - P$ so the code works wherever it is mapped.
Table-based relocations write the offset of the symbol slot so the text stays shareable while only the table holds absolute addresses.
The complexity is $O(M + R)$, where $M$ is the total input size for merging and $R$ is the relocation count, plus symbol-table hashing costs that behave like linear time in practice.
Relaxation and identical-code folding add passes proportional to the section sizes they rewrite.
Overflow checks are not optional decoration: a 32-bit field holding a 64-bit address must fail loudly rather than truncate silently.
% TODO(cite): need a primary reference for dynamic-loader symbol binding.

Position independence changes which fixups the static linker leaves behind.
A \term{position-independent executable} uses PC-relative and table-based addressing so the loader can map it at any base.
A \term{global offset table} holds absolute addresses of shared data, while a \term{procedure linkage table} routes function calls through lazily bound stubs.
The static linker builds those tables and emits the dynamic relocations the loader will consume.
Lazy binding resolves each function on first call through a resolver rather than at startup, which shortens launch time at the price of a first-call indirection.
Security features such as read-only-after-relocation ordering, immediate binding options, and control-flow annotations ride on the same tables.

\section{Worked Example}\label{sec:ch29-worked}

We link two tiny objects by hand with concrete numbers.
Object \code{main.o} has a 64-byte text section starting at placeholder zero and one data word.
Object \code{util.o} has a 32-byte text section and one data word.
The link script places text at \code{0x1000} with 16-byte alignment and data at the next 16-byte boundary after text.
Layout: \code{main.o} text occupies \code{0x1000} through \code{0x103F}, \code{util.o} text occupies \code{0x1040} through \code{0x105F}, and data starts at \code{0x1060} with the \code{util.o} counter at \code{0x1060} and the \code{main.o} word at \code{0x1068} after padding.
Symbol \code{helper} resolves to \code{0x1040}.
Symbol \code{counter} resolves to \code{0x1060}.

Now patch three relocations step by step.
First, a direct call at \code{0x1010} to \code{helper} with PC-relative 32 computation and addend \code{-4}: the patched value is \code{0x1040 - 4 - 0x1010 = 0x2C}, so the call encodes offset \code{0x2C}.
Check the arithmetic: target \code{0x1040}, place \code{0x1010}, processor adds the instruction length before branching, hence the \code{-4} addend.
Second, an absolute 64-bit pointer to \code{counter} in data at \code{0x1068} writes \code{0x1060} plus its addend of \code{0} with no overflow because the field is 64 bits wide.
Third, a data reference to \code{helper} through the table writes the table slot offset rather than \code{0x1040} directly, keeping the referencing section shareable.
Verify each result: the call lands on the first byte of \code{helper}, the pointer loads the counter address, and the table slot will hold \code{0x1040} after load-time fixup.
A single off-by-alignment error anywhere in this chain would shift every later address and corrupt all three patches.

Loading continues where linking stopped for dynamic builds.
The loader maps the text segment read-execute, the read-only segment read-only, and the data segment read-write at the chosen base.
It processes dynamic relocations in dependency order, resolves lazy function slots on first call, runs initializers, and then jumps to the entry point.
Unloading reverses the order with finalizers.
Startup latency depends on relocation count, so table design and binding policy are performance features, not just correctness machinery.

\section{Implementation Notes}\label{sec:ch29-implementation}

Production linkers optimize for throughput on huge inputs: parallel section hashing, concurrent relocation patching, and incremental modes for edit-build-debug cycles.
LLVM lld, GNU ld, and mold differ in parallelism strategy and script compatibility but agree on the skeleton above.
Dead-section stripping uses reachability from entry points and exported symbols; identical-code folding hashes instruction bytes and merges duplicates with careful exception-table handling.
Relaxation rewrites conservative long sequences into short ones once distances are final, for example shrinking a global access into a PC-relative load.
Each of these changes bytes after addresses were assigned, so the linker iterates layout and patching to a fixed point with bounded passes.

Dynamic loading adds the runtime half of the contract.
The loader enforces segment permissions, applies address-space randomization offsets, binds symbols breadth-first across dependencies, and honors visibility and versioning rules.
Common failures include missing shared objects on the search path, version mismatches that resolve the wrong definition, and initialization-order cycles across libraries.
Diagnose them with loader tracing and symbol lookup logs before suspecting code generation.
Reproducible builds add deterministic section ordering, stable hash seeds, and normalized timestamps so identical sources produce identical images.

Testing this layer means comparing addresses, not just outputs.
Check map files for section placement, verify relocations with object-dump tools, and run the image under multiple base addresses to prove position independence.
Combine objects built at different optimization levels to catch convention drift.
Useful compiler-side context for object emission and symbol handling appears in~\cite{appel-1998-modern} and see~\cite{morgan-1998-optimizing}.
Validation thinking for preserving behavior across link-time transformations follows see~\cite{safonov-2010-trustworthy}.
Rules for linkage, storage duration, and observable effects at the language level are defined in~\cite{wg14-2011-c11}.

\section{Exercises}\label{sec:ch29-exercises}

\begin{enumerate}
\item Two objects define text sizes 48 and 64 bytes with text base \code{0x2000} and 16-byte alignment. \label{ex:ch29-layout} Compute the start address of each section and the first data address, showing padding explicitly.
\item A PC-relative call sits at \code{0x3000} targeting \code{0x3080} with addend \code{-4}. \label{ex:ch29-pcrel} Compute the encoded offset by hand and verify that it lands on the target.
\item An archive provides \code{foo} and \code{bar} in one member, but the program uses only \code{foo}. \label{ex:ch29-archive} Explain what the linker extracts, what stays out, and how link order could break the build.
\end{enumerate}

\section{Further Reading}\label{sec:ch29-reading}

Object emission, symbol tables, and backend-to-linker interfaces appear in~\cite{appel-1998-modern}.
Link-time optimization, section handling, and target concerns are developed in~\cite{morgan-1998-optimizing}.
Validation of whole-program transformations spanning compile and link time is discussed in~\cite{safonov-2010-trustworthy}.
Linkage, storage, and conformance rules at the language level appear in~\cite{wg14-2011-c11}.
